\section{The Science and technology Question}

\begin{quotex}
The next question from the interview published in ``La Nation Europeenne'' deals with science and technology.
\end{quotex}

\textbf{Q}. Your judgment of science and technology seems, in your work, to be negative. What are the reasons for your position? Don't you believe that material conquests and the elimination of hunger and poverty will lead to the tackling of spiritual problems with more energy?

\textbf{A}. As to the second point you brought up, I will say that, as there exists a state of degradation due to poverty, so there exists a state of degradation due to affluence and prosperity. The ``Welfare State'' in which on can no longer speak of hunger and poverty, are far from engendering an increase in true spirituality; on the contrary, there is established a violent and destructive form of revolt by the new generations against the system in its entirety and against an existence devoid of any meaning (USA, England, Scandinavia). The problem consists, instead, in setting a just limit, constraining the frenzy of a capitalistic economy, the creator of artificial needs, and freeing the individual from his growing dependence on social and productive mechanisms. It would be necessary to establish a balance. Just until recently, the Japanese have provided the example of a balance of this type; it modernized itself and didn't leave itself behind the West in the scientific and technical domains, while safeguarding its specific traditions. But today the situation is quite different.

There is another fundamental point to emphasize: it is difficult to adopt science and technology, while circumscribing them within the limits of a civilization's material means and instruments, that is, while preserving, in regards to them, a certain distance; on the contrary, it is practicably inevitable the it becomes impregnated with the world conceptions on which modern profane science bases itself, conceptions that are practically inculcated in our spirits by the methods of the customary methods of instruction and that has, on the spiritual plane, a destructive effect. The very concept of true knowledge is thus totally distorted.

\begin{quotex}
This speaks for itself. If not, see the Sorcerer's Apprentice.

However, if science is merely instrumental, then what is the nature of ``true knowledge''. Can it, too, be ``inculcated''?
\end{quotex}

\flrightit{Posted on 2008-01-06 by Cologero}

\begin{center}* * *\end{center}

\begin{footnotesize}
\begin{sffamily}
\texttt{CaseyAnn on 2013-07-03 at 14:06 said:}

Excellent. This is always a struggle to explain to progressive secularists who see themselves as more charitable because they believe in an abstract brotherhood-of-all-nations void of any meaning.

\hfill

\texttt{Mark Citadel on 2015-02-08 at 15:17 said:}

The best reason to be very wary about the goals of transhumanism, an element prominent within NeoReaction is pointed out by Evola. I would like to add another.

``Simply put, the technology that helps us survive has never kept up with the technology designed to kill us''

\url{http://citadelfoundations.blogspot.com/2015/02/the-utopian-scientists.html}



\hfill

\end{sffamily}
\end{footnotesize}

